\documentclass[11pt,a4paper]{article}
\usepackage[T1]{fontenc}
\usepackage{lmodern,amsmath,amssymb,booktabs,microtype}
\usepackage[margin=24mm]{geometry}
\usepackage[colorlinks=true,urlcolor=blue,linkcolor=blue]{hyperref}
\newcommand{\EE}{\mathbb E}
\newcommand{\tr}{\operatorname{tr}}
\newcommand{\PP}{\mathsf P_h}
\newcommand{\FF}{\mathsf F}
\title{Finite-time memory from two returns\\A controlled reduction to ordinary charge propagation}
\author{Alan Day \and Chengheng\thanks{AI research assistant (OpenAI).}}
\date{Local research note --- 6 November 2026}
\begin{document}\maketitle
\begin{abstract}
For the common-coupling complex SYK charge-removal dilation, we
reduce the actual finite-memory covariance to an auxiliary kernel
with ordinary Hamiltonian propagation. A finite Volterra iteration
and uniform Gaussian force moments give a two-return approximation
with collective-field error $O(p^{-2})$ and absolute covariance error
$O(p^{-1})$. Time exponentials are kept intact. The remaining task
is the limit of the resulting pinched, common-coupling products.
\end{abstract}
\section{State, operators and the result}
Use standard CAR, $\hbar=1$, string-time units, and
$H=\sum_{I,J}J_{IJ}c_I^\dagger c_J$, where $I,J$ are increasing
$p$-subsets, $c_J=(c_J^\dagger)^\dagger$, $J=J^\dagger$ is Gaussian, and
\begin{equation}
 \EE J_{IJ}J_{KL}=w_N\delta_{IL}\delta_{JK},\qquad
 w_N=4^p\frac{J_s^2}{\lambda_N}\left(\frac{p!}{N^p}\right)^2,
 \quad \lambda_N=p^2/N.
\end{equation}
Take even $N\to\infty$, $\lambda_N\in[\lambda_-,\lambda_+]\subset(0,\infty)$,
with $J_s$ fixed. The annealed state is the normalized full trace
at $\beta=\mu=0$: sector weights are $b_m=2^{-N}\binom Nm$.
Write $h=H_m$, $d_m=\binom Nm$, $\tr_m=\operatorname{Tr}/d_m$. For $m>0$ set
\begin{align}
 V\psi&=m^{-1/2}\sum_i|i\rangle c_i\psi,&
 K&=1_N\otimes H_{m-1},& Q&=1-VV^\dagger,\\
 a&=V^\dagger KV=(1-p/m)h,& R&=QKV,& D&=QKQ,\\
 G&=R^\dagger R,& M(s)&=R^\dagger e^{iDs}R,&
 \Lambda(s)&=R^\dagger e^{iKs}R.
\end{align}
All sectors use the same $J$. Let $\PP X=\sum_eP_eXP_e$, including
whole degenerate energy blocks. The actual dynamic field is
\begin{equation}
 \FF_m(s)=(1-\PP)\big[m e^{-ihs/2}M(s)e^{-ihs/2}\big].
\end{equation}
Inactive $m<p$ sectors, including $m=0$, are zero. Define convolution
$A*B(s)=\int_0^s A(s-t)B(t)\,dt$, $U(s)=e^{ias}$ and
$J_n=\Lambda*(U*\Lambda)^{*n}$, with $J_0=\Lambda$.
Replacing $M$ by $J_0+J_1+J_2$ defines $\FF^{[2]}$.
We prove, uniformly for $|s|,|u|\leq S$ and all real $T$,
\begin{align}
 \|\FF(s)-\FF^{[2]}(s)\|_{2,\mathrm{ann}}&=O_S(p^{-2}),\label{eq:mainfield}\\
 |C_N(s,u;T)-C_N^{[2]}(s,u;T)|&=O_S(p^{-1}),\label{eq:maincov}
\end{align}
where $C_N=\sum_m b_m\EE\tr_m\FF_m(s)^\dagger
e^{iTh}\FF_m(u)e^{-iTh}$ and $C_N^{[2]}$ uses the proxy fields.

\newpage
\section{Exact finite iteration, not a time Taylor truncation}
Removing projected propagation by auxiliary kernels is an established
Mori method \cite{MR,WZS}. Here finite matrices allow a direct proof.
In $P\oplus Q$, identify $P$ with the source space using $V$:
\begin{equation}
 K=\begin{pmatrix}a&R^\dagger\\R&D|_Q\end{pmatrix}.
\end{equation}
The $Q$-block Duhamel equations give two transitions and hence $i^2=-1$:
\begin{equation}
 \Lambda=M-M*U*\Lambda,\qquad M=\Lambda+M*U*\Lambda.\label{eq:volterra}
\end{equation}
For example, eliminating the $P$-block of $e^{iKs}R$ yields a
convolution of $e^{iDs}$, $R$, $U$, $R^\dagger$, and the full
$Q$-block. Multiplying by $R^\dagger$ and $R$ gives
\eqref{eq:volterra}, preserving the matrix order.
Associativity gives, for any fixed integer $j\geq0$,
\begin{equation}
 M=\sum_{n=0}^jJ_n+\mathcal E_j,\qquad
 \mathcal E_j=M*(U*\Lambda)^{*(j+1)}.\label{eq:finite}
\end{equation}
This is a finite identity with an exact remainder. No convergence
of an infinite Neumann or Taylor series is assumed.

Normalize every rectangular Schatten trace by the \emph{same}
source dimension $d_m$. An integrand in $\mathcal E_j$ contains
$2j+4$ factors $R$ or $R^\dagger$, separated by unitary
propagators. Schatten H\"older with exponent $4j+8$ bounds its
Hilbert--Schmidt norm by $(\tr_mG^{2j+4})^{1/2}$.
The $2j+2$ integrations have simplex volume
$|s|^{2j+2}/(2j+2)!$. Therefore, samplewise,
\begin{align}
 \|\mathcal E_j(s)\|_{2,m}
 &\leq\frac{|s|^{2j+2}}{(2j+2)!}(\tr_mG^{2j+4})^{1/2},\label{eq:rem}\\
 \|J_n(s)\|_{2,m}
 &\leq\frac{|s|^{2n}}{(2n)!}(\tr_mG^{2n+2})^{1/2},\label{eq:term}\\
 \|M(s)\|_{2,m}&\leq(\tr_mG^2)^{1/2}.
\end{align}
For negative times use adjoints: reversing the simplex order gives
$J_n(-s)=J_n(s)^\dagger$. Equivalently the oriented convolutions
contain an even number of integrations. Thus all bounds apply to
both time signs.

An elementary sign check is $K=\sigma_x$, $V=(1,0)^t$:
$a=0$, $R=(0,1)^t$, $M=1$. The first terms are
\begin{equation}
 J_0=\cos s,\qquad J_1=\frac{s\sin s}{2},\qquad
 J_2=\frac{s\sin s-s^2\cos s}{8}.
\end{equation}
In particular $M-\Lambda=s^2G^2/2+O(s^3)$ for general matrices.
The finite-return expressions themselves retain all time orders.

\newpage
\section{Force square functions and uniform high moments}
Let $F_i=[H,c_i]$, $\mathcal A=\sum_iF_i^\dagger F_i$.
The exact dilation identity already used in the project is
\begin{equation}
 G=\mathcal A_m/m-(p/m)^2h^2,\qquad 0\leq G\leq\mathcal A_m/m.
\end{equation}
Here a direct Gaussian bound replaces any need to count high-order
marked chord diagrams. For $T_{IJ}=c_I^\dagger c_J$,
$[T_{IJ},c_i]$ vanishes unless $i\in I$; otherwise it is
$\pm c_{I\setminus i}^\dagger c_J$.
The covariance maps of $F_i$ are diagonal in occupation states.
At charge $m$ and occupation $\eta=n_i$, their eigenvalues divided
by $w_N$ are
\begin{align}
 f_{\rm in}(m,0)&=\binom mp\binom{N-m+p-1}{p-1},\\
 f_{\rm in}(m,1)&=\binom{m-1}p\binom{N-m+p}{p-1}
 +\binom{m-1}{p-1}\binom{N-m+p-1}{p-1},\\
 f_{\rm out}(m,\eta)&=\binom{m-\eta}{p-1}\binom{N-m+p-1}p.
\end{align}
Invalid binomials are zero. Count the annihilated $J$ first for
$f_{\rm in}$, separating whether it includes $i$; count the occupied
$I\setminus i$ first for $f_{\rm out}$. Overlaps $I\cap J$ are retained.
Both counts are at most
$(N+p)^{2p-1}/[4^{p-1}p!(p-1)!]$: bound binomials by powers,
pair $p-1$ powers with $xy\leq(x+y)^2/4$, and bound the remaining
power by $N+p$. Hence
\begin{equation}
 \EE F_i^\dagger F_i,\ \EE F_iF_i^\dagger\ \leq v_N1,\quad
 v_N=\frac{4J_s^2}{\lambda_N}\frac pN(1+p/N)^{2p-1}
 \leq\frac{4J_s^2}{\lambda_N}\frac pN e^{2\lambda_N}.\label{eq:v}
\end{equation}

For clarity, the Gaussian matrix moment step is also explicit.
If $X=\sum_\alpha g_\alpha B_\alpha=X^\dagger$ and
$\sum B_\alpha^2\leq v1$, Gaussian integration by parts gives
\begin{equation}
 \EE\operatorname{Tr}X^{2k}
 =\sum_{\ell=0}^{2k-2}\EE\sum_\alpha
 \operatorname{Tr}B_\alpha X^\ell B_\alpha X^{2k-2-\ell}
 \leq(2k-1)v\EE\operatorname{Tr}X^{2k-2}.
\end{equation}
To see the inequality, diagonalize $X$, take absolute values and
apply weighted arithmetic--geometric mean to
$|x_b|^\ell|x_a|^{2k-2-\ell}$. Symmetry of
$|(B_\alpha)_{ab}|^2$ gives
$\operatorname{Tr}(\sum B_\alpha^2)|X|^{2k-2}$.
Apply the recurrence to the Hermitian dilation of the rectangular
$F_i:\mathcal H_m\to\mathcal H_{m-1}$, stopping at its second moment.
The factor two from the dilation cancels, so
\begin{equation}
 \EE\tr_m(F_i^\dagger F_i)^k\leq(2k-1)!!v_N^k,
 \qquad \EE\tr_m\mathcal A_m^k\leq(2k-1)!!(Nv_N)^k.\label{eq:amom}
\end{equation}
The last step is Minkowski in the annealed Schatten $k$ norm;
independence between the different $F_i$ is unnecessary.

\newpage
\section{Charge summation and the collective covariance error}
Equations \eqref{eq:v}--\eqref{eq:amom} hold uniformly in charge.
For $k\geq2$, eigenvalue monotonicity of $0\leq G\leq\mathcal A_m/m$
gives $\tr_mG^k\leq m^{-k}\tr_m\mathcal A_m^k$; this does not
assert that the power map is operator monotone. Eventually $p\leq N/4$.
Split the binomial sum at $m=N/4$ and use
$\Pr(m<N/4)\leq e^{-N/8}$. Since active sectors have $m\geq p$,
\begin{align}
 B_{k,N}&:=\sum_m b_m\EE m^2\tr_mG^k\\
 &\leq(2k-1)!!(Nv_N)^k
 \left[(N/4)^{2-k}+p^{2-k}e^{-N/8}\right]
 =O_k(p^{4-k}).\label{eq:weighted}
\end{align}
There is no central-window limit exchange in this argument.
Only $k=2,4,6,8$ are needed below. Exterior centering is unitary
on both sides and $1-\PP$ is an orthogonal Hilbert--Schmidt
projection. Thus, denoting the centered, pinched $mJ_n$ by $X_n$,
\begin{align}
 \|\FF(s)\|_{2,\mathrm{ann}}&\leq\sqrt{B_{2,N}}=O(p),\\
 \|X_n(s)\|_{2,\mathrm{ann}}&\leq
 \frac{S^{2n}}{(2n)!}\sqrt{B_{2n+2,N}}=O_S(p^{1-n}),\\
 \|\FF(s)-\FF^{[j]}(s)\|_{2,\mathrm{ann}}&\leq
 \epsilon_{j,N}:=\frac{S^{2j+2}}{(2j+2)!}\sqrt{B_{2j+4,N}}
 =O_S(p^{-j}).\label{eq:scaled}
\end{align}
Here $\FF^{[j]}=\sum_{n=0}^jX_n$. The exact finite-$N$ energy
pinching is retained in every term.

Conjugation by $e^{iTh}$ is an isometry, including for sample-dependent
fields. Cauchy--Schwarz on the direct sum and Gaussian probability
space gives
\begin{equation}
 |C_N-C_N^{[j]}|\leq
 2\sqrt{B_{2,N}}\,\epsilon_{j,N}+\epsilon_{j,N}^2.\label{eq:covbound}
\end{equation}
At $j=1$, the field error is $O(p^{-1})$, but this only proves
an $O(1)$ covariance error using the available field bound.
At $j=2$, \eqref{eq:covbound} is $O(p^{-1})$, proving
\eqref{eq:mainfield}--\eqref{eq:maincov}. An $O(1)$ bound on
the true finite-time field was not presumed.

There is also a six-product form of the remaining calculation.
Define
\[
 \langle X,Y\rangle_T=\sum_m b_m\EE\tr_m
 X_m^\dagger e^{iTh}Y_me^{-iTh}.
\]
Then
\begin{equation}
 C_N(s,u;T)=
 \sum_{\substack{n,l\geq0\\n+l\leq2}}
 \langle X_n(s),X_l(u)\rangle_T+O_S(p^{-1}).\label{eq:six}
\end{equation}
The omitted pairs $(1,2),(2,1),(2,2)$ are bounded by
$O(p^{-1}),O(p^{-1}),O(p^{-2})$. For numerical positivity,
the full nine-product $C_N^{[2]}$ remains preferable.
All estimates fix $S$ and $\lambda_->0$; they do not cover
growing memory times or a simultaneous $\lambda\to0$ limit.

\newpage
\section{Ordinary propagation and a stable spectral formula}
Let $F(s)=V^\dagger e^{iKs}V$ (this $F$ is a compressed propagator).
Expanding $R=KV-Va$ gives the exact identity
\begin{equation}
 \Lambda(s)=-F''(s)+i\{a,F'(s)\}+aF(s)a.\label{eq:derivative}
\end{equation}
The sign of the last term follows directly from this expansion;
it should not be copied from a different Mori projection convention.
Keep $a=(1-p/m)h$ exactly.

Choose eigenbases $h|a\rangle=E_a|a\rangle$ and
$H_{m-1}|\alpha\rangle=\varepsilon_\alpha|\alpha\rangle$.
Put $a_a=(1-p/m)E_a$, $c^i_{\alpha a}=\langle\alpha|c_i|a\rangle$, and
\begin{equation}
 W^\alpha_{ab}=\frac1m\sum_i(c^i_{\alpha a})^*c^i_{\alpha b}
 (\varepsilon_\alpha-a_a)(\varepsilon_\alpha-a_b).
\end{equation}
Then $\Lambda_{ab}(s)=\sum_\alpha W^\alpha_{ab}e^{i\varepsilon_\alpha s}$.
For $r\geq0$, define the time-simplex integral
\begin{equation}
 \mathcal I_r(s;\xi_0,\ldots,\xi_r)=
 \int_{\substack{t_j\geq0\\\sum t_j=s}}
 e^{i\sum_j\xi_jt_j}\,dt_0\cdots dt_{r-1}
 =i^{-r}[e^{is\,\cdot}]^{[r]}(\xi_0,\ldots,\xi_r),\label{eq:simplex}
\end{equation}
initially for $s\geq0$ and then continued to real $s$.
The first two returns are the finite spectral paths
\begin{align}
 (J_1)_{ab}&=\sum_{\alpha,\beta,c}W^\alpha_{ac}W^\beta_{cb}
 \mathcal I_2(s;\varepsilon_\alpha,a_c,\varepsilon_\beta),\\
 (J_2)_{ab}&=\sum_{\alpha,\beta,\gamma,c,d}
 W^\alpha_{ac}W^\beta_{cd}W^\gamma_{db}
 \mathcal I_4(s;\varepsilon_\alpha,a_c,\varepsilon_\beta,a_d,
 \varepsilon_\gamma).
\end{align}
Coincident energies use confluent divided differences.
One stable evaluation is the top-right entry of $e^{sB}$,
where $B$ has diagonal $i\xi_j$ and unit superdiagonal. It
automatically gives $s^re^{i\xi s}/r!$ when all nodes coincide.

Set $L=J_0+J_1+J_2$. For one sector and one realization the
proxy covariance is
\begin{equation}
 \frac{m^2}{d_m}\sum_{E_a\ne E_b}
 e^{iT(E_a-E_b)+i(s-u)(E_a+E_b)/2}
 \overline{L_{ab}(s)}L_{ab}(u).\label{eq:phase}
\end{equation}
Summing with $b_m$ and averaging gives $C_N^{[2]}$. Equal-energy
blocks, not merely diagonal entries, are omitted. The phase in
\eqref{eq:phase} keeps both the center time and the two memory times.
No independent-copy replacement of $J$, annealed spectral factorization,
or replacement of the exact pinching has been made.

\newpage
\section{Checks and the next calculation}
The new driver \texttt{memory\_return\_expansion.py} in \texttt{code/}
stores \texttt{finite-checks.json} in the stage results directory,
without overwriting existing evidence. It checks:
\begin{itemize}
 \item Five complete CAR systems through $N=5$, including $p=3$:
 864 force-covariance diagonal entries agree exactly as integers;
 off-diagonal covariance entries vanish. Another 118,132 rational
 cases verify the combinatorial upper bound for $2\leq N\leq70$.
 \item Four common-coupling systems $(N,p,m)=(3,1,2)$, $(4,2,3)$,
 $(5,2,3)$, $(6,2,3)$ at $s=-0.8,0.35,1.1,2$:
 all 48 finite-return remainder bounds hold. The derivative identity,
 negative-time adjoints and covariance phases are independently checked.
 \item A block-triangular exponential computes the return convolutions;
 a separate spectral-path sum checks $J_0,J_1,J_2$ at $(3,1,2)$.
 Its largest absolute-scaled matrix discrepancy is $9.39\times10^{-16}$.
 Five scalar cases check the explicit trigonometric formulas, including
 the confluent zero-energy return nodes.
\end{itemize}
The Hamiltonians in the finite propagation tests are seeded Gaussian
integer Hermitian coefficient matrices, divided by their operator
norms. Their seed, exact coefficient hash and scale are saved.
They test identities, not ensemble asymptotics. Arithmetic is
\texttt{complex128}; $2\times10^{-12}$ identity tolerances and
$2\times10^{-13}$ bound allowances are empirical roundoff tolerances.

At $s=2$, the observed remainder/bound ratios are:
\begin{center}
\begin{tabular}{rrrrr}\toprule
$N$&$p,m$&zero returns&one return&two returns\\\midrule
3&1,2&0.852425&0.877750&0.899381\\
4&2,3&0.783725&0.838456&0.871689\\
5&2,3&0.913060&0.934683&0.947976\\
6&2,3&0.899579&0.921893&0.936414\\\bottomrule
\end{tabular}
\end{center}
The $(3,1,2)$ dynamic pinched covariance is numerically zero;
the other three cases supply nonzero checks. For $(4,2,3)$ at
$(s,u,T)=(-0.8,1.1,0.67)$, the two-return covariance error is
$1.71146\times10^{-5}$, below the direct Cauchy--Schwarz bound
$2.47819\times10^{-5}$. These finite errors are retained.

\paragraph{Next step.}
Evaluate the six common-coupling, pinched products in
\eqref{eq:six}, establish finite-$s$ tightness and identify their
limit. This is a concrete route from the microscopic bilocal to
the Susskind-line bulk meson problem. The present reduction does
not yet determine that limit or a bulk mass/space dictionary.
The earlier three Taylor-coefficient spectra remain unchanged.

\begin{thebibliography}{9}
\bibitem{MR} A.~Montoya-Castillo and D.~R.~Reichman,
\emph{Approximate but Accurate Quantum Dynamics from the Mori
Formalism: I. Nonequilibrium Dynamics},
\href{https://arxiv.org/abs/1603.01903}{arXiv:1603.01903},
Sec.~III, Eqs.~(11)--(26). Auxiliary-kernel removal and derivative
rewritings are established methods; that paper credits Shi--Geva.
\bibitem{WZS} C.~Widder, J.~Zimmer and T.~Schilling,
\emph{On the generalized Langevin equation and the Mori projection
operator technique}, \href{https://arxiv.org/abs/2503.20457}{arXiv:2503.20457},
Secs.~2--3. General context for projected evolution and Volterra
equations; no unbounded-operator theorem is required here.
\end{thebibliography}
\end{document}
